\section{Scope and conditional statement}

Let \(X\) be a connected compact K\"ahler manifold. Write
\(X^\beth\) for its MP ideal spectrum and \(o\) for the zero
semivaluation. Let \(\alpha\in H^{1,1}(X,\mathbf R)\) be nef, let
\(\theta\) be a smooth closed representative of \(\alpha\), and let
\(\omega\) be a K\"ahler form, with \(\beta=\{\omega\}\).
Nefness implies that \(\alpha+\varepsilon\beta\) is K\"ahler for
every \(\varepsilon>0\); it does not imply that the particular form
\(\theta+\varepsilon\omega\) is positive.

Denote the full DXZ space, with the inverse-system definition of
\cite[Definition~3.5]{XiaOps}, by \(\mathcal N(\theta)\). Its
components are the normalized \(\mathcal I\)-model test curves of
\cite[Definition~3.2]{XiaOps}, including the prime version in that
definition. We do not assume positive mass for an arbitrary component.
For \(\Gamma\in\mathcal N(\theta)\), write \(c_\Gamma\in\mathbf R\)
for its common endpoint. Its raw transform is
\begin{equation}\label{eq:p143n-raw}
 F_\Gamma(v)=\sup_{\tau<c_\Gamma}
 \left\{\tau-\sup_{m\geq1}\frac{1}{m}
 v\!\left(\mathcal I(m\Gamma^{\theta+\eta}_\tau)\right)\right\},
 \qquad v\in X^\beth,
\end{equation}
where \(\eta\) is any K\"ahler form. Transitions preserve the
multiplier ideals, so this is independent of \(\eta\). We use extended
real arithmetic with \(\tau-\infty=-\infty\). In particular,
\(F_\Gamma(o)=c_\Gamma\) and \(F_\Gamma\leq c_\Gamma\).

The sole new mathematical dependency of this packet is the following
conditional comparison, already isolated in a separate packet.
\begin{assumption}\label{hyp:p143n-kahler}
For every K\"ahler class \(\gamma\) on \(X\), the raw transform is
an order isomorphism
\[
 \mathcal N(\gamma)\longrightarrow\PSH_{\mathrm{MP}}(\gamma)
\]
as actual functions on all of \(X^\beth\). Here the DXZ space for a
class uses its standard change-of-representative identification.
\end{assumption}
This input is the full K\"ahler comparison, not just its finite-energy
part and not just equality on divisorial points. Its raw-USC and
dominated-divisorial dependencies, and the comparison itself, remain
subject to the separate project review gates.

Define a set of actual functions, not a quotient by divisorial equality,
by
\begin{equation}\label{eq:p143n-intersection}
 \mathcal Q_\omega(\alpha)
 :=\bigcap_{\varepsilon>0}
 \PSH_{\mathrm{MP}}(\alpha+\varepsilon\beta).
\end{equation}

\begin{theorem}\label{thm:p143n-nef}
Assume \cref{hyp:p143n-kahler}. The raw transform induces an order
isomorphism
\[
 \mathcal N(\theta)\longrightarrow\mathcal Q_\omega(\alpha),
 \qquad \Gamma\longmapsto F_\Gamma.
\]
The right-hand side is independent of the auxiliary K\"ahler form
\(\omega\), and the isomorphism is unchanged by the standard DXZ
identification for a different smooth representative of \(\alpha\).
The statement includes nef non-big classes, with no positive-volume or
finite-energy assumption on \(\Gamma\).
\end{theorem}

We do not identify \(\mathcal Q_\omega(\alpha)\) with the decreasing
completion of relatively nef models in the fixed boundary class
\(\alpha\). Nor does the theorem assert a comparison for a non-nef
pseudoeffective class, an unnormalized zero-volume energy or measure
theory, or retirement of the full P14.3 question.
